\scanpage{73}
% lines 1–3, opening identification and formula
identifiant $E$ à $\deleted{\unsure{E}}L_1\otimes L_2\oplus\underline O_X$ grâce au scindage \uline{$x_1*x_2$}, on trouve, par un calcul immédiat$^{(*)}$
\[
T=\deleted{\unsure{\chi(x_1*x_2)}}\unsure{\chi|\gamma}(x_1*x_2)\mathbin{\unsure{-|+}}\deleted{\unsure{\chi}}b_1\otimes b_2
\]
% lines 3–7, two annotations below the identification
\margin{* NB\\ $\pi_T(x_1*x_2)=\chi\,x_1*x_2-T$ \deleted{\unsure{$\chi$}}\\ donc $T=\chi\,x_1\otimes x_2-\pi_T(x_1*x_2)$\\ Or $\pi_T(x_1*x_2)=\pi_{T_1}(\deleted{\unsure{u}}x_1)\otimes\pi_{T_2}(x_2)$\\ $\forall\,x_1\in E_1,\ x_2\in E_2$}
\margin{Re NB : on \unsure{vérifie|vérifierait} aussitôt que ça ne change \unsure{pas} si on remplace $x_1$ par $x_1+u_1$ ou $x_2$ par $x_2+u_2$ ($u_1\in\Gamma L'_1,\ u_2\in\Gamma L_2$)\\ Attention signe !}
% lines 8–12, three braced definitions
Pour éliminer le signe, il faudrait poser
\[
\left\{\begin{aligned}
b_{T_1,x_1}&=-b_1&\quad\text{i.e.}\quad b_{T_1,x_1}&=\chi x_1-T_1=\pi_{T_1}(x_1)\\
b_{T_2,x_2}&=-b_2&\text{i.e.}\quad b_{T_2,x_2}&=\chi x_2-T_2=\pi_{T_2}(x_{\unsure{1|2}})\\
b_{T,x_1*x_2}&=-b&\text{i.e.}\qquad b&=\chi\,x_1\mathbin{\unsure{*|\otimes}}x_2-T=\pi_T(x)\\
&&&(x=x_1\otimes x_2)
\end{aligned}\right.
\]
% lines 13–17, consequence and formula for T
et on trouvera
\[
\deleted{\unsure{b_2|b}}\ b=-b_{T,x_1*x_2}=-b_1\otimes b_2=-b_{T_1,x_1}\otimes b_{T_2,x_2}
\]
d'où
\[
\boxed{b_{T,x_1*x_2}=b_{T_1,x_1}\otimes b_{T_2,x_2}}
\]

On trouve alors
\[
T=\chi\cdot(x_1*x_2)-\underbrace{b_{T,x_1*x_2}}_{b_{T_1,x_1}\otimes b_{T_{\unsure{1|2}},x_2}}=
\]
\centerline{---}
% lines 18–25, formulary
\uline{Formulaire récapitulatif} pour ce formalisme
\[
(x_1,x_2)\longrightarrow x_1*x_2:\quad E_1\times E_2\longrightarrow E_1*E_2,
\]
où $E_i$ \inserted{$(i=1,2)$} est l'espace affine enveloppant d'un torseur $P_i$ sous un module \deleted{$\unsure{M_i}$} $L_i$, torseur muni d'un hom de $\chi$-trivialisation \inserted{($\chi\in\Gamma(X,O)$, donné)}
\[
\pi_i:P_i\longrightarrow L_i
\]
satisfaisant (par définition de ``$\chi$ trivialisation'')
\[
\pi_i(x_i+u_i)=\pi_i(x_i)+\chi\,u_i\qquad x_i\in\Gamma P_i,\ u_i\in\Gamma L_i
\]
% archive sheet number: 72
